% !PRINTED 308
\weberpage{308}

%<<CONT
not torn apart, but only interchanged among one another. Moreover, in two of these pairs the elements are interchanged each time. If we therefore form a function like
\weq{9}{V = (v_\infty - v_0)(v_1 - v_4)(v_2 - v_3),}
then it remains unchanged under $B, C, U$, and hence under the whole system (3), while by repeated application of the cyclic permutation $(0, 1, 2, 3, 4)$, that is, of the substitution
\[
A = \begin{pmatrix} 1, & 0 \\ 1, & 1 \end{pmatrix}
\]
it takes \gesperrt{five} different values. But the whole group $\mathfrak{L}_0$ can be composed from $A$ and $B$ [§ 79, (10)], and $V$ therefore takes no more than five values under the application of $\mathfrak{L}_0$.

The five values of $V$ are the roots of an equation of the fifth degree.

For $p = 7$ we have
\[
B = \left(z, -\frac{1}{z}\right), \quad C = (z, 4z), \quad U = \left(z, \frac{\mp 2 + z}{3 \mp z}\right).
\]

For a better survey we also represent $B, C, U$ by the \gesperrt{cyclic permutations} which they produce among the eight values of $z$:
\[
\begin{aligned}
B &= (0, \infty)(1, -1)(2, 3)(-2, -3), \\
C &= (1, -3, 2)(-1, 3, -2), \\
U &= (\infty, \mp 1, \pm 1, \pm 3)(0, \mp 3, \mp 2, \pm 2),
\end{aligned}
\]
and from this one obtains the \gesperrt{seven-valued function}
\weq{10}{V = (v_\infty - v_0)(v_{\mp 1} - v_{\mp 3})(v_{\pm 1} - v_{\mp 2})(v_{\pm 3} - v_{\pm 2}),}
in which the one or the other sign may be taken.

For $p = 11$ we have
\[
B = \left(z, -\frac{1}{z}\right), \quad C = (z, 4z), \quad U = \left(z, \frac{\mp 1 + z}{2 \mp z}\right),
\]
or, represented by the cycles:
\[
\begin{aligned}
B &= (0, \infty)(1, -1)(2, 5)(-2, -5)(3, -4)(-3, 4), \\
C &= (1, 4, 5, -2, 3)(-1, -4, -5, 2, -3), \\
U &= (\infty, \mp 1, \pm 3, \mp 2, \pm 2)(0, \pm 5, \mp 5, \mp 4, \pm 1),
\end{aligned}
\]
from which the two 11-valued functions
\weq{11}{V = (v_\infty - v_0)(v_{\mp 1} - v_{\pm 5})(v_{\pm 3} - v_{\mp 5})(v_{\mp 2} - v_{\mp 4})(v_{\pm 2} - v_{\pm 1})(v_{\pm 4} - v_{\mp 3}).}
